\section{C++ 面试代码质量和容器工程细节}

算法面试中的 C++ 代码不只是能 AC，还要让状态、所有权和复杂度可读。函数参数只读大对象应使用 const 引用，返回结果时依赖移动优化即可；循环中遍历复杂对象使用 const auto\&，不要无意复制字符串、向量或节点记录。局部常量要有名字，尤其是方向数组、无穷大、字符状态、取模常数和哨兵值。

整数类型选择要主动。数组下标和 size 返回 size\_t，但很多题需要与 -1 哨兵或差值比较，直接混用无符号可能产生隐蔽 bug。常见做法是在函数开始把 size 转成 int，并保证输入规模在 int 范围内；涉及总和、乘积、面积、距离时使用 \code{std::int64_t}。二分平方、堆距离、图最短路都要优先防溢出。

容器 API 语义要讲清。unordered\_map 的 operator[] 在 key 不存在时会插入默认值，适合计数累加，不适合只检查存在。find 返回迭代器，适合查询后读取值。contains 只判断存在，如果随后还要取值会导致二次查找。vector 的 reserve 只改容量，resize 才改大小。priority\_queue 的 top 只保证极值，不能遍历出有序序列。

迭代器失效是工程题常见追问。vector 扩容后所有指针和迭代器失效；erase 中间元素会让后续位置移动；unordered\_map rehash 会让迭代器失效；list 的 splice 能移动节点且保持其他迭代器稳定。LRU 缓存为什么用 list 加哈希表，核心就是哈希定位和链表节点稳定的组合。

代码组织要让不变量有位置。二分答案把检查函数单独写出来，图最短路把松弛逻辑集中，回溯把 push、搜索、pop 成对出现，DP 把初始化和主转移分开。这样做能让读者看到状态生命周期。把所有逻辑塞进一个巨大循环，短期可能省行数，长期会让边界和证明都变模糊。

面试时可以主动说明替代方案。第 K 大可以排序、堆、快速选择；岛屿数量可以 DFS、BFS、并查集；最小体力路径可以 Dijkstra、二分加 BFS、并查集排序边；单词接龙可以普通 BFS、通配桶、双向 BFS。比较方案不是炫耀，而是说明你知道当前选择依赖什么约束。

Linux 源码阅读放到后面的独立专题中集中学习。面试代码部分只需要先记住一个工程判断：容器选择不只看复杂度，还要看对象生命周期、内存分配、并发保护和数据是否需要被多个索引同时组织。

实验建议：写一个前缀和计数题，分别用 contains 加 operator[]、find、直接 operator[] 三种方式实现，观察语义差异。写一个 vector 扩容实验，保存元素地址后继续 push\_back，验证地址失效。写一个 LRU，先用 vector 保存顺序，再用 list 加 unordered\_map，比较移动到头部的成本。

\begin{principlebox}{本节复盘}
阅读本节后，请回到相邻章节的例题，把暴力瓶颈、优化依据、状态不变量、C++ 容器成本和边界测试重新写一遍。能把同一思想迁移到变体题，才算真正掌握。
\end{principlebox}
